\section{A population information reduction and a radial obstruction}
\label{sec:population-information}
\noindent\textbf{PROVED-ON-PAPER.} These statements concern the actual iid
hidden-gap experiment. They neither reveal a sampled gap nor provide an
oracle for the unknown sign allocation. The reduction below is finite-prefix;
the stated testing consequence uses almost-sure termination separately at
the two endpoint laws. No unproved differentiability claim at an unbounded
stopping time is used. No Lean formalization or priority claim is made.

\subsection{Directional information inequalities tensorize over interleaving}
Fix a known magnitude law on $(0,a]$, $a<1/2$, and a path
$q_t=q_*+tu$, $t\in[\alpha,\beta]$, of measurable feasible sign
allocations in $\mathcal Q_\rho$. The algorithm is independent of $t$.
When $\alpha<\beta$, feasibility at the two endpoints implies
$|u|\le1/(\beta-\alpha)$ almost everywhere; the degenerate path is trivial.
For a finite bit word $h$ write $f_t(h)$ for its probability on one
fresh coin, and set
\[
 s_t(h)=\partial_t\log f_t(h),\quad
 p_t(b\mid h)=f_t(hb)/f_t(h),\quad
 v_t(h)=\sum_{b=0}^1p_t(b\mid h)
                  [s_t(hb)-s_t(h)]^2.
\]
Here $f_t(h)$ is affine in $t$ and strictly positive, because all Bernoulli
parameters lie in $[1/2-a,1/2+a]$. A bounded number of observations therefore
has a well-defined ordinary Fisher information. Include the independent
random seed in each transcript; its law is fixed when taking a derivative.

\begin{lemma}[Finite-prefix information/work direct sum]
\label{pop:direct-sum}
Suppose a measurable function $\kappa(t)\ge0$ has this property: for
each fixed $t$, every randomized single-coin policy with a deterministic
finite cap on its number of flips satisfies
\[
 I_t(\text{seed and stopped coin history})
 =\mathbb E_t\sum_{k<\tau}v_t(H_k)
 \le\kappa(t)\mathbb E_t\tau.
\]
Then every population algorithm with at most $N<\infty$ total flips,
arbitrary interleaving and adaptive creation of iid coins satisfies
\[
 I_t(\text{full population transcript})
 =\mathbb E_t\sum_{n<T}v_t(H_{i_n,n})
 \le\kappa(t)\mathbb E_tT.
\]
Here $H_{i_n,n}$ is the previous local history of the coin selected for
the next observation. The single-coin premise may grant its policy knowledge
of the reference value $t$; the derivative keeps that policy fixed.
\end{lemma}
\begin{proof}
Conditioning on the seed, each possible transcript has likelihood equal to
the product of $f_t(h_i)$ over the local histories of its distinct coins.
The adaptive scheduling and stopping choices are already determined by
the displayed seed and past observations, so contribute no likelihood
factor depending on $t$. Revealing the next bit changes the score by
$s_t(hb)-s_t(h)$, which has conditional mean zero. Squaring the finite
sum of score increments and taking expectations makes the cross terms
vanish. This proves the two information identities.

For the inequality fix a reference $t$ and a creation index $i\le N$,
where coins are indexed by their first charged observation. Run a copy of
the capped population policy with its $i$th such coin supplied by one
genuine coin, and simulate all other coins using independent tapes drawn
from the known reference law $q_t$. The simulated tapes are part of an
independent seed. They are held fixed in law while differentiating the
genuine coin's distribution. If the $i$th coin is never requested, take
zero genuine flips. This is a valid one-coin policy using at most $N$
genuine flips; simulated observations and arithmetic are only auxiliary
randomization in this information comparison. Its expected information
and work at the reference law are exactly the contributions of coin $i$
to the two displayed population sums. The hypothesis bounds that
contribution by $\kappa(t)\mathbb E_t T_i$. Summing over $i\le N$
proves the result. No simulated reference coins are supplied to the
original statistical algorithm.
\end{proof}

The proof is not specific to Bernoulli observations: it applies to any
independent reservoir observation channel having positive differentiable
finite-history likelihoods, provided differentiation under the seed and
observation integrals is justified and the stated one-coin information/work
premise holds. This is a conditional transfer principle, not a guessed
universal information functional for other channels or cost models.

\begin{corollary}[Expected-budget testing consequence]
\label{pop:testing}
Let an unrestricted population procedure satisfy
$\sup_{t\in[\alpha,\beta]}\mathbb E_tT\le B<\infty$.
If the premise of Lemma~\ref{pop:direct-sum} holds, then the total variation
distance between its two terminal transcript laws is at most
\[
 \frac{\sqrt B}{2}\int_\alpha^\beta\sqrt{\kappa(t)}\,dt.
\]
Consequently, if this quantity is at most $1/2$ and
$\Delta=|\theta(q_\beta)-\theta(q_\alpha)|$, its worst-case RMS error
over these two laws is at least $\Delta/4$.
\end{corollary}
\begin{proof}
For every finite prefix capped at $N$ flips, integrate its score derivative
along the path and apply Cauchy--Schwarz:
\[
 \|P_{\beta,N}-P_{\alpha,N}\|_{\rm TV}
 \le\tfrac12\int_\alpha^\beta\sqrt{I_t(P_{t,N})}\,dt
 \le\tfrac{\sqrt B}{2}\int_\alpha^\beta\sqrt{\kappa(t)}\,dt.
\]
Finite prefix densities are finite sums over observation/action branches
conditional on the seed, so differentiating is legitimate; the same
identity extends by integration over the fixed independent seed.
Finite expected work implies almost-sure termination separately at each
endpoint. For any terminal event $E$, its intersection with $\{T\le N\}$
is determined by that prefix. Thus its endpoint probability difference
is bounded by the preceding total-variation bound plus
$\mathbb P_\alpha(T>N)+\mathbb P_\beta(T>N)$, which tends to zero.
Take the supremum over $E$. Finally threshold an estimator at the midpoint
of the two target values. The sum of its two testing errors is at least
$1-\operatorname{TV}\ge1/2$, whereas Markov's inequality bounds that
sum by $8R^2/\Delta^2$. Hence $R\ge\Delta/4$.
\end{proof}

The missing hypothesis is an upper bound on \emph{directional Fisher
information per expected flip for every single-coin policy}. A lower bound
on the uniform variance--work product of estimators of the entire
functional does not provide this hypothesis. The next example shows why
the most natural radial path cannot supply the extra critical logarithm.

\subsection{The radial path has too much information for that converse}
Fix again $a=1/4$, the exact quadratic law, and define $b\in(0,a)$ by
\[
 w_b(d)=\min\{1,b^2/d^2\},\qquad
 \mathbb E w_b(D)=\rho.
\]
Set $\ell=\log(a/b)$ and $L=1+\log(1/\rho)$. Then
$\rho=(b/a)^2(1+2\ell)$ and
$L=1+2\ell-\log(1+2\ell)$.
Consider the genuine one-dimensional submodel
\[
 q_\eta(d)=\tfrac12(1+\eta)w_b(d),\qquad -1/2\le\eta\le1/2.
\]
It belongs to $\mathcal Q_\rho$, and its target derivative is $\rho/2$.

\begin{proposition}\label{pop:radial-information}
At $\eta=0$, the supremum of single-coin information/work ratios over
finite-cap policies is at least
\[
 \frac{a^2\rho}{32L}.
\]
Thus no uniform upper bound of order $a^2\rho/[L(1+\log L)]$ can hold
for this directional information envelope as $\rho\downarrow0$.
\end{proposition}
\begin{proof}
Let $M=\lceil8a/b\rceil$. On the genuine coin, run the integer walk
with increments $2X-1$ from zero until it hits $-1$ or $M$; output the
indicator $A$ of hitting $M$. Let $\pi_\pm(d)$ be its success probabilities
conditional on mean $1/2\pm d$. The elementary gambler's-ruin formula gives
\[
 \pi_+(d)=\frac{1-r}{1-r^{M+1}},\quad
 r=\frac{1-2d}{1+2d},\qquad
 \pi_-(d)=\frac{R-1}{R^{M+1}-1},\quad R=r^{-1}.
\]
These follow by solving the two-boundary recurrence, or by summing its
geometric first differences. In particular
\begin{equation}\label{pop:ruin-bounds}
 \frac83d\le\pi_+(d)\le4d+\frac1{M+1},\qquad
 \pi_-(d)\le\frac{8d}{e^{4(M+1)d}-1}.
\end{equation}
For the upper positive bound, write
$\pi_+=(1-r)+r^{M+1}\pi_+$ and use
$r^{M+1}(M+1)\le\sum_{k=0}^M r^k=1/\pi_+$.
The negative bound uses $R-1\le8d$ and $\log R\ge4d$.

Write $F=\mathbb E\pi_-(D)$ and
$A_+=\mathbb E[q_0(D)\pi_+(D)]$. Expanding
$(e^x-1)^{-1}=\sum_{k\ge1}e^{-kx}$ and using
$\sum_{k\ge1}k^{-3}<2$ yields
\[
 F\le\frac{16}{a^2}\int_0^\infty
           \frac{d^2}{e^{4(M+1)d}-1}\,dd
 <\frac1{a^2(M+1)^3}
 \le\frac{b^3}{512a^5}.
\]
Direct integration gives
$\mathbb E[D w_b(D)]=2b^2/a-4b^3/(3a^2)\ge2b^2/(3a)$.
Hence \eqref{pop:ruin-bounds} gives
\[
 A_+\ge\tfrac43\mathbb E[D w_b(D)]\ge\frac89\frac{b^2}{a},
 \qquad \frac F{A_+}\le\frac{9b}{4096a^4}
 \le\frac9{64}<\frac14.
\]
Let $z_\eta=\mathbb P_\eta(A=1)$. Because this single-coin law is affine
in $\eta$,
\[
 z_0\le A_++F\le\tfrac54 A_+,\qquad
 z'_0=\mathbb E[q_0(\pi_+-\pi_-)]\ge A_+-F\ge\tfrac34 A_+.
\]
Both boundary events have strictly positive probability at every signed
gap; hence $0<z_0<1$.
The information in the output indicator alone therefore satisfies
\begin{equation}\label{pop:radial-I}
 I_0(A)=\frac{(z'_0)^2}{z_0(1-z_0)}
 \ge\frac9{20}A_+\ge\frac25\frac{b^2}{a}.
\end{equation}

For fixed signed $d$, the hitting time has finite mean: in every block of
$M+1$ observations there is a positive, uniform-in-the-current-state chance
of hitting a boundary, because both bit probabilities are at least $1/4$.
The stopped walk identity (or the same finite difference recurrence) gives
\[
 \mathbb E_{-d}\tau=\frac{1-(M+1)\pi_-(d)}{2d}\le\frac1{2d},
 \qquad
 \mathbb E_{+d}\tau
 =\frac{(M+1)\pi_+(d)-1}{2d}
 \le2(M+1)+\frac1{2d}.
\]
At $\eta=0$, integration and
$\mathbb E[w_b(D)/D]\le4b/a^2$ imply
\begin{align*}
 \mathbb E_0\tau
 &\le\frac1a+(M+1)\rho+\frac b{a^2}\\
 &\le8+8(b/a)(1+2\ell)+2\rho<25.
\end{align*}
The last inequality uses $e^{-\ell}(1+2\ell)\le2$ for $\ell\ge0$,
which follows from $e^\ell\ge1+\ell$.
Now set $u=1+2\ell$. The elementary inequality $\log u\le u/2$ gives
$L=u-\log u\ge u/2$, so $b^2\ge a^2\rho/(2L)$.
Combining this with \eqref{pop:radial-I} and $a=1/4$ shows
\[
 \frac{I_0(A)}{\mathbb E_0\tau}
 \ge\frac{2b^2}{125a}
 \ge\frac{a\rho}{125L}
 =\frac{a^2\rho}{(125/4)L}
 \ge\frac{a^2\rho}{32L}.
\]
To obtain the asserted finite-cap result, stop after at most $N$ flips
and output one only if the upper boundary was already hit. For $N\ge M$
this event and its complement both have positive probability. Its success
probability and its $\eta$ derivative converge to $z_0,z'_0$ by dominated
convergence: the mixture probabilities are affine in $\eta$ and
$|q'_\eta|=q_0\le1/2$. Its work increases to $\mathbb E_0\tau$.
The finite-cap output information therefore converges to $I_0(A)$, and
the full transcript has at least this information. Taking the supremum
over finite caps proves the proposition.
\end{proof}

\subsection{An estimator at the candidate order on the known radial submodel}
\begin{corollary}\label{pop:radial-upper}
Under the known radial submodel of Proposition~\ref{pop:radial-information},
and the ideal arithmetic convention allowing exact evaluation of known
one-dimensional calibration integrals, there is an observable estimator
with uniform expected work at most $B>0$ and worst-case RMS at most
\[
 \min\left\{\rho/4,\sqrt{160\rho L/B}\right\}.
\]
The guarantee is only over $q_\eta$, $|\eta|\le1/2$, and not over
arbitrary sign allocations in $\mathcal Q_\rho$.
\end{corollary}
\begin{proof}
All expectations in the following calibration constants are integrals
against the known exact quadratic law. Set
\[
 d_0=\mathbb E[q_0(D)(\pi_+(D)-\pi_-(D))],\qquad
 F=\mathbb E\pi_-(D),\qquad
 Y=\frac{\rho}{2d_0}(A-F).
\]
These constants depend only on the known submodel, not on $\eta$ or on
the realized gap. The previous proof gives $d_0\ge(2/3)b^2/a>0$ and
$F\le d_0/3$. Since $z_\eta=F+(1+\eta)d_0$,
$\mathbb E_\eta Y=(1+\eta)\rho/2=\theta(q_\eta)$, and
\[
 \operatorname{Var}_\eta(Y)
 \le\frac{\rho^2}{4d_0^2}z_\eta
 \le\frac{11\rho^2}{24d_0}
 \le\frac{11a\rho^2}{16b^2}
 \le\frac{11}{8a}\rho L=\frac{11}{2}\rho L.
\]
The upper bound for $\pi_+$ in \eqref{pop:ruin-bounds}, retaining the
subtraction of one in the work formula, improves the positive conditional
work bound to $\mathbb E_{+d}\tau\le2(M+1)$. Thus uniformly in this
submodel,
\[
 \mathbb E_\eta\tau
 \le\frac1a+2(M+1)\theta(q_\eta)
 \le4+\frac32(M+1)\rho
 \le4+\frac32(16+1/2)<29.
\]
Each test terminates almost surely by the finite-interval argument above.

Independently draw $N\sim\operatorname{Poisson}(B/29)$ and run $N$
independent copies on genuinely fresh coins, with independent seeds. Put
$\lambda=B/29$ and output
\[
 \widehat\theta=\rho/2+\lambda^{-1}
                   \sum_{i=1}^N(Y_i-\rho/2).
\]
This is a finite sum almost surely and costs at most $29\mathbb EN=B$
in expectation, uniformly in $\eta$. The compound-Poisson identities
(obtained directly by conditioning on $N$) give
$\mathbb E\widehat\theta=\theta$ and
\[
 \operatorname{Var}(\widehat\theta)
 =\lambda^{-1}\mathbb E_\eta(Y-\rho/2)^2
 \le\frac{29}{B}\left(\frac{11}{2}\rho L+\frac{\rho^2}{16}\right)
 <\frac{160\rho L}{B},
\]
using $\rho^2/16\le\rho L/64$.
The deterministic output $\rho/2$ has worst-case error $\rho/4$ on this
submodel. Choose the smaller stated guarantee using known $B,\rho$.
\end{proof}

\paragraph{Calibration and practical scope.}
This is a proved estimator in the stated ideal computation model, not a
held-out benchmark or claim of efficient numerical calibration. Its
one-dimensional integrands are explicit and continuous, using the fair
limit $\pi_\pm(0)=1/(M+1)$, and the cutoff $b$ is the unique root of a
known strictly increasing continuous function. For a fixed justified integer
$M$ and a fixed known radial profile, these calibration integrals can be
approximated to arbitrary certified precision. The following error bound
holds with $b,M$ fixed and approximates only $F,d_0$; it does not certify
root-finding for $b$ or rounding $\lceil8a/b\rceil$. No such executed
certification is claimed. For explicit error accounting, if
$|\widehat d_0-d_0|\le d_0/2$ and
$\widehat Y=\rho(A-\widehat F)/(2\widehat d_0)$, then pointwise
\[
 |\widehat Y-Y|
 \le\rho\left(\frac{|\widehat d_0-d_0|}{d_0^2}
                 +\frac{|\widehat F-F|}{d_0}\right)=:\zeta.
\]
Using the same independent Poisson count, the RMS difference of the
calibrated and ideal final outputs is at most
$\zeta\sqrt{1+1/\lambda}$, because it is pointwise at most
$\zeta N/\lambda$. Add this term to the ideal RMS guarantee; flip costs
are unchanged. No actual numerical integration or implementation run is
claimed here, and any computational calibration cost must be reported
separately in a practical comparison.

\paragraph{Pseudocode.}
Compute $b,M,F,d_0$ from the known radial model and select the better of
the trivial radius and the displayed Poisson bound. If selecting the
trivial bound, output $\rho/2$ immediately. Otherwise sample the independent
Poisson count. For each of its genuinely fresh coins, accumulate $2X-1$
until the running sum hits $-1$ or $M$, record the upper-hit indicator,
and convert it to $Y$ by the displayed calibration. Return the centered
Poisson sum. No posterior involving the unknown $\eta$, no true gap,
and no selected-magnitude request is used.

The proposition disproves a particular proposed directional information
upper bound, not the unrestricted minimax rate. The calibrated estimator
knows the radial profile $w_b$ through the specified submodel and estimates
$\int q\,d\mu$ only on that submodel. It charges every
genuine flip and never uses the realized magnitude. The full nonparametric
problem could still require an entropy loss arising from multiple nuisance
directions, or could admit a nonlinear population construction absent from
the single-coin product argument.
